¿Cómo puede un procesador el doble de rápido perder? Por lo que vimos al principio: con contención, lo caro no es calcular, sino mover la línea compartida de un núcleo a otro.

En el Ultra 9 medimos esos viajes: entre 35 y 86 nanosegundos según dónde estén los núcleos. Y sin fijar los hilos, Linux los reparte entre núcleos P y E, que están lejos.

El i5 es un solo bloque de seis núcleos, y cada núcleo ejecuta dos hilos que comparten sus cachés. Si dos hilos caen en el mismo núcleo, la línea ni siquiera viaja.

Tenemos una pista fuerte en el propio Ultra 9: dos hilos que comparten caché L2 rinden casi el doble que cualquier otra pareja. Es la misma física.

Pero en el i5 no lo medimos, así que es una hipótesis. Lo decimos claro.
